% =====================================================================
%  coq_coverage.tex
%
%  A reader's map from the paper's theory (theory/paper/article.tex and
%  its appendices) to the machine-checked Coq development in
%  proof_verification/coq-formal/.
%
%  Standalone document: compile with
%      pdflatex coq_coverage.tex   (run twice for cross-references)
%  or  latexmk -pdf coq_coverage.tex
%  Uses only standard TeX Live packages (amsmath, amssymb, booktabs,
%  longtable, listings, hyperref, geometry, xcolor).
% =====================================================================
\documentclass[11pt]{article}

\usepackage[a4paper,margin=1in]{geometry}
\usepackage{amsmath,amssymb,amsthm}
\usepackage{booktabs}
\usepackage{longtable}
\usepackage{array}
\usepackage{xcolor}
\usepackage{listings}
\usepackage[hidelinks]{hyperref}
\usepackage{enumitem}

% ---- Coq identifier / label helpers -------------------------------------
% \detokenize renders underscores literally, so Coq names need no escaping.
\newcommand{\coq}[1]{\texttt{\detokenize{#1}}}
\newcommand{\lbl}[1]{\texttt{\detokenize{#1}}}
\newcommand{\file}[1]{\texttt{\detokenize{#1}}}

% Long typewriter identifiers/paths are unbreakable; let TeX break the line
% before them rather than run them into the margin.
\emergencystretch=3em
\sloppy

% ---- a lightweight Coq listing style ------------------------------------
\definecolor{coqkw}{rgb}{0.30,0.12,0.55}
\definecolor{coqcmt}{rgb}{0.30,0.45,0.30}
\lstdefinelanguage{Coq}{
  morekeywords={Theorem,Lemma,Corollary,Definition,Record,Proof,Qed,forall,exists,
    Section,Variable,Variables,Hypothesis,Hypotheses,Require,Import,From,fun,
    intros,unfold,rewrite,apply,field,split,nra,lra,psatz,pose,assert},
  sensitive=true,
  morecomment=[n]{(*}{*)},
}
\lstset{
  language=Coq,
  basicstyle=\ttfamily\small,
  keywordstyle=\color{coqkw}\bfseries,
  commentstyle=\color{coqcmt}\itshape,
  columns=fullflexible,
  keepspaces=true,
  breaklines=true,
  frame=single,
  framesep=4pt,
  xleftmargin=2pt,
}

% ---- math shorthands matching the paper ---------------------------------
\newcommand{\sigt}{\widetilde{\sigma}}
\newcommand{\Cinv}{C_{\mathrm{inv}}}
\newcommand{\aK}{\alpha_K}
\newcommand{\Reh}{\mathrm{Re}_h}
\newcommand{\Dah}{\mathrm{Da}_h}
\newcommand{\Cstab}{C_{\mathrm{stab}}}
\renewcommand{\arraystretch}{1.25}

\newcolumntype{L}[1]{>{\raggedright\arraybackslash}p{#1}}

\title{\bfseries Machine-checked coverage map:\\
       from the paper's theory to the Coq formalisation}
\author{Coverage note accompanying \file{proof_verification/coq-formal/}}
\date{Revised 2026-07-11 \quad (12 files, $\approx$384 theorems;
      Rocq/Coq build verified: version 9.1.1)}

\begin{document}
\maketitle

\begin{abstract}
\noindent
This note is a reader's map between two artefacts in this repository: the
mathematical theory of the paper \emph{``A stabilized finite element method
for incompressible, inertial flows in inhomogeneous porous media''}
(Casas, Gonz\'alez-Us\'ua, Codina, de-Pouplana), whose \LaTeX{} sources live
under \file{theory/paper/}, and the Coq formalisation under
\file{proof_verification/coq-formal/}, which re-proves the paper's algebraic
and analytic \emph{claims} from the axioms of the real numbers and has them
re-checked by Coq's trusted kernel. \textbf{The headline of the current
revision: \lbl{lemma:Stability} and \lbl{lemma:Continuity} are now complete
machine-checked theorems} (\coq{abstract_stability},
\coq{abstract_continuity_sharp}, \coq{abstract_continuity}), proved over an
abstract inner-product / mesh interface from a short, \emph{named,
quantitative} trusted base of standard analytic facts. Section~\ref{sec:coq}
is a primer on how Coq (now \emph{Rocq}) works and what ``machine-checked''
buys. Section~\ref{sec:map} explains the translation convention --- including
how a lemma with functional-analytic hypotheses is honestly encoded.
Section~\ref{sec:crosswalk} is the crosswalk itself, one table per Coq file.
Section~\ref{sec:left} states precisely what is \emph{not} inside the kernel
and by what other means (a hand audit, the SymPy suite, a Lean~4 roadmap, or
the convergence tests) each residual item is addressed.
Section~\ref{sec:amend} records the theory-side amendments applied to keep the
manuscript $100\%$ compatible with what was verified.
\end{abstract}

\tableofcontents

% =====================================================================
\section{Purpose and scope}\label{sec:purpose}
% =====================================================================

The Coq development mirrors, file by file, the SymPy ``verification scripts''
suite in \file{proof_verification/sympy/}: where SymPy checks an identity
by symbolic computation, the Coq files \emph{prove} it. The suite is
\textbf{12 files, $\approx 384$ theorems, all proved} --- there is no
\coq{Admitted}, no \coq{Axiom} and no \coq{admit} anywhere in the sources
(\coq{grep} confirms this). It was developed against Coq~8.18 using
\emph{only} the standard library, and it compiles and kernel-checks unchanged
under the current Rocq~9.1.1.

What changed since the previous revision is a step up in \emph{kind}, not just
in count. Previously the Coq layer certified the \emph{algebraic skeleton} of
the two central lemmas (parameter identities, coefficient expansions, the jump
algebra) while the surrounding \emph{functional-analytic} argument (testing
with a discrete field, integration by parts, inverse estimates, summation over
the mesh) was only hand-audited. Four new files
(\file{AbstractSums.v}, \file{InnerSpace.v}, \file{AbstractStability.v},
\file{AbstractContinuity.v}) close that gap: they prove \lbl{lemma:Stability}
and \lbl{lemma:Continuity} \emph{as complete theorems} over an abstract
inner-product / mesh interface, reducing each to a short list of
explicitly-named standard analytic facts (Green-type identities and
inverse-type estimates) that appear as \emph{hypotheses} of the abstract
theorem. The elemental and discrete Cauchy--Schwarz inequalities, the Young
step, the difference-of-squares cancellation, all eighteen term bounds, the
jump treatment, and the assembly over elements and faces are now inside the
kernel.

Two boundaries still frame everything.
\begin{itemize}[nosep]
  \item \textbf{Coverage boundary.} Coq certifies the mathematics down to a
    named trusted base. That base --- the divergence-theorem / integration-by-parts
    identities and the finite-element inverse estimates --- is textbook material
    (Brenner--Scott, Ern--Guermond), hand-audited against the manuscript and
    slated for a Lean~4 formalisation
    (Section~\ref{sec:left}). It is \emph{stated quantitatively} as hypotheses,
    so the abstract theorems are genuine implications, not sketches. Only the
    interpolation-theoretic replacements of \lbl{lem:continterp} remain
    hand-audited without a Coq statement.
  \item \textbf{Meaning boundary.} Coq verifies the \emph{mathematics as
    stated}: that each identity, inequality or limit follows from the
    real-number axioms and the named base. It does \emph{not} certify that the
    statements model the physics, nor that the Julia/Kratos solver implements
    them --- that is the job of the method-of-manufactured-solutions
    convergence tests elsewhere in the codebase.
\end{itemize}

% =====================================================================
\section{A short primer: how Coq works and what ``verified'' means}\label{sec:coq}
% =====================================================================

\subsection{Proof assistant, kernel, and the de Bruijn criterion}

Coq (renamed \emph{Rocq} from version~9) is an interactive proof assistant
built on a small, fixed logic --- the Calculus of Inductive Constructions. A
Coq \emph{proof} is not prose: it is a formal term whose \emph{type} is the
statement being proved (the ``propositions-as-types'' correspondence). The
user constructs that term interactively with \emph{tactics}, but tactics are
untrusted conveniences; what matters is the object they build.

The trust rests on the \textbf{de Bruijn criterion}: a large, human-friendly
front end may do arbitrary heuristics, but the final proof term is re-examined
by a \emph{small, fixed kernel} that only knows the typing rules of the logic.
If the kernel accepts the term, the theorem holds. In this development the
kernel is exercised twice:
\begin{enumerate}[nosep]
  \item \coq{coqc} type-checks each file as it compiles it;
  \item \coq{coqchk} --- a \emph{separate}, independent re-implementation of
    the kernel --- re-verifies the compiled objects (\coq{.vo} files) from
    scratch, with no trust in \coq{coqc}. The build script runs it and reports
    \emph{``coqchk: all modules kernel-verified''}.
\end{enumerate}

\subsection{The trust base --- and why ``no \texttt{Admitted}, no \texttt{Axiom}'' matters}

Believing a Coq theorem means trusting exactly three things: the logic itself
(the Calculus of Inductive Constructions, studied for decades), the kernel
implementation, and the \emph{statement} --- one must still read what was
proved, including its hypotheses. Everything else (tactics, notations, the
elaborator) cannot make a false theorem be accepted.

Two escape hatches would silently widen that base, and both are audited out
here:
\begin{itemize}[nosep]
  \item \coq{Admitted} / \coq{admit} closes a goal \emph{without} a proof; a
    single one would let anything ``follow''.
  \item \coq{Axiom} postulates a statement as true without proof.
\end{itemize}
Neither appears anywhere in the twelve files. Consequently every result below
is reduced to the standard-library axiomatisation of the real numbers (plus,
for the two abstract lemmas, the explicitly-listed analytic hypotheses ---
Section~\ref{sec:left}). The suite uses \emph{no external libraries} (no
Mathcomp, no Coquelicot): the analysis is done on \coq{Coq.Reals} with the
arithmetic tactics \coq{lra}/\coq{nra} (linear and nonlinear real arithmetic)
and \coq{field} (rational-function identities). Those tactics are
decision/normalisation procedures whose \emph{output} is a kernel-checked
term --- using them does not weaken the guarantee.

\subsection{The abstract-interface method (new in this revision)}

The two main lemmas are proved not against a concrete function space but
against a \emph{minimal algebraic interface} that the concrete objects satisfy.
Two devices carry this.

\emph{A pre-Hilbert record.} \file{InnerSpace.v} models the elemental
$L^2(K)$ pairings by a four-field record: a carrier with addition and scalar
multiplication and a form \coq{ip} that is symmetric, additive and homogeneous
in its first argument, and positive-semidefinite.

\begin{lstlisting}
Record PreHilbert := mkPreHilbert {
  carrier  :> Type;
  vadd     : carrier -> carrier -> carrier;
  vscal    : R -> carrier -> carrier;
  ip       : carrier -> carrier -> R;
  ip_sym   : forall x y, ip x y = ip y x;
  ip_add_l : forall x y z, ip (vadd x y) z = ip x z + ip y z;
  ip_scal_l: forall a x y, ip (vscal a x) y = a * ip x y;
  ip_pos   : forall x, 0 <= ip x x }.
\end{lstlisting}

\noindent
From these four axioms the file \emph{derives} (does not assume) full
bilinearity, the Cauchy--Schwarz inequality (by the discriminant argument,
with the degenerate direction handled), the norm expansions of sums and
differences, the triangle inequality, the parametrized Young inequality at
vector level, and the difference-of-squares identity
$\langle -a+b,\,a+b\rangle = \lVert b\rVert^2 - \lVert a\rVert^2$ that drives
the ASGS cancellation. Any concrete weighted $L^2(K)$ is an instance, so the
proofs transfer for free.

\emph{A named trusted base.} A lemma such as \lbl{lemma:Stability} needs one
genuinely analytic input --- a weighted inverse estimate --- that stdlib Coq
cannot prove without finite-element scaling theory. Rather than pretend
otherwise, the estimate is written as an explicit, \emph{quantitative}
\coq{Hypothesis} of the theorem. The theorem is then a real implication:
\emph{if} the inverse estimate holds with the stated constant, \emph{then} the
stability bound holds with the stated constant. This is the honest encoding of
a ``cited-standard'' fact, and it makes the residual trusted base a finite,
readable list (Section~\ref{sec:left}) rather than a gap.

\subsection{Reading Coq syntax}\label{sec:syntax}

Table~\ref{tab:syntax} is enough to read every entry in the crosswalk.

\begin{table}[h]
\centering
\small
\begin{tabular}{@{}L{0.34\textwidth}L{0.58\textwidth}@{}}
\toprule
\textbf{Coq} & \textbf{Meaning} \\
\midrule
\coq{Definition n : R := e.}       & Names a real-valued quantity $n := e$. \\
\coq{Record R := mk { f : T; ax : P }.} & A structure bundling fields \coq{f} with axioms \coq{ax} they satisfy. \\
\coq{Theorem}/\coq{Lemma}/\coq{Corollary n : P.} & Asserts $P$; the following \coq{Proof. ... Qed.} is its (kernel-checked) proof. \\
\coq{Section S. ... End S.}        & A scope sharing common parameters and hypotheses. \\
\coq{Variables (x y : R).}         & Introduces symbolic reals $x,y$ for the section. \\
\coq{Hypothesis h : 0 < x.}        & A standing assumption; every theorem in the section is conditioned on it. \\
\coq{forall x, ...} / \coq{exists x, ...} & Universal / existential quantifier. \\
\coq{<< x , y >>}                  & The pre-Hilbert inner product \coq{ip x y}; \coq{nrm x} $= \sqrt{\langle x,x\rangle}$. \\
\coq{Rsum Th f}                    & The finite sum $\sum_{k\in\mathtt{Th}} f(k)$ over the mesh element (or face) list. \\
\coq{a <= b}, \coq{a < b}, \coq{a = b} & Order and equality on $\mathbb{R}$. \\
\coq{sqrt}, \coq{Rabs}, \coq{exp}, \coq{sin} & Standard-library real functions. \\
\coq{derivable_pt_lim f x d}       & ``$f$ has derivative $d$ at $x$'' (the stdlib derivative predicate). \\
\bottomrule
\end{tabular}
\caption{The Coq vocabulary used in the statements.}
\label{tab:syntax}
\end{table}

A note on total division. In \coq{Coq.Reals}, $x/0$ is defined and equals $0$;
the development never relies on this --- every division either is harmless by
construction (its statement carries no hypothesis and the denominator is a
declared positive quantity) or is guarded by an explicit non-vanishing
hypothesis.

% =====================================================================
\section{How paper statements map to Coq}\label{sec:map}
% =====================================================================

\subsection{Standing hypotheses become section variables}

The paper carries global positivity conventions --- mesh size $h>0$, viscosity
$\nu>0$, elemental porosity $\aK>0$ (with $\aK\le 1$ in the appendix),
$\sigma\ge 0$, non-dimensional numbers $\Reh,\Dah\ge 0$, constants
$c_1,c_2,\Cinv,\xi>0$, profile $0<\alpha_0<1$. In Coq these are the
\coq{Hypothesis} lines of each \coq{Section}; every theorem there is read as
``under these positivity assumptions, \dots''.

\subsection{Naming dictionary}

The Coq identifiers transliterate the paper's symbols. The recurring ones:

\begin{center}
\small
\begin{tabular}{@{}llll@{}}
\toprule
Paper & Coq & Paper & Coq \\
\midrule
$\tau_1$ (\lbl{eq:Tau1Final})       & \coq{tau1}       & $\varphi_1=\aK\,\tau_{1,\mathrm{NS}}^{-1}$ & \coq{phi1} \\
$\tau_2$ (\lbl{eq:Tau2Final})       & \coq{tau2}       & $\tau_{1,\mathrm{NS}}$ (\lbl{eq:TauNavierStokes}) & \coq{tau1NS} \\
$\sigt$ (\lbl{eq:SigmaAlpha})       & \coq{sigt}/\coq{sg} & $\bar C_{\mathrm{inv}}$ (\lbl{lem:winv}) & \coq{Cb} \\
$\aK$                               & \coq{alphaK}/\coq{aK}     & $|\mathbf a|_{\infty,K}$   & \coq{am} \\
$\Cinv$ (\lbl{eq:inverse})          & \coq{Cinv}       & $\xi$ (Young parameter)   & \coq{xi} \\
$\Reh,\Dah$                         & \coq{Re},\coq{Da} & $\varepsilon$ (mass penalty) & \coq{eps} \\
$B_S(U_h,U_h)$                      & \coq{BS}         & $|||U_h|||^2$ (\lbl{eq:triplenorm}) & \coq{NormSq} \\
\bottomrule
\end{tabular}
\end{center}

\subsection{A worked example, read side by side}

The sharpest scalar statement of the stability section is the \emph{velocity
slack identity}: the gap between the velocity coefficient of the final
stability estimate (\lbl{eq:VelocityCoefficientBound}) and $C_u\,\sigt$ is not
merely non-negative --- it equals a specific product \emph{exactly}. In Coq:

\begin{lstlisting}
Definition C_u : R := 1 - xi * Cinv^2 / c1.

Theorem velocity_slack_identity :
  u_855 - C_u * sigt
    = alphaK * tau1 * sigma * (xi * Cinv^2 / c1) * (c2 * amag / h).

Corollary velocity_coefficient_lower_bound : u_855 >= C_u * sigt.
\end{lstlisting}

\noindent
Read: \emph{under the section hypotheses}, the difference of the two
real-valued expressions equals $\aK\tau_1\sigma\,(\xi\Cinv^2/c_1)(c_2|\mathbf
a|/h)\ge 0$; hence the lower bound. The proof rewrites with the identity
$\sigt=\sigma\varphi_1\tau_1$ and discharges the rational-function equality
with \coq{field}.

\subsection{How a lemma with analytic hypotheses maps}\label{sec:absexample}

The abstract theorems read the same way, with the trusted-base facts as named
hypotheses. \lbl{lemma:Stability} is \coq{abstract_stability}: under the
section's positivity plus $c_1>2\xi\bar C_{\mathrm{inv}}^2$, $\xi>2$,
$C_2<1$, and the three analytic hypotheses

\begin{lstlisting}
(* HBS: Galerkin testing + adjoint expansion (eq:StabilityEstimate) *)
Hypothesis HBS :
  BS = 2*nu * Rsum Th (fun k => << gu k , gu k >>)
     + sigma * Rsum Th (fun k => << uu k , uu k >>)
     + eps   * Rsum Th (fun k => << pp k , pp k >>)
     - Rsum Th (fun k => t1 k * << Lstar_m k , L_m k >>)
     - Rsum Th (fun k => t2 k * << Lstar_c k , L_c k >>).

(* S3: the weighted inverse estimate eq:winv-divvisc (lem:winv) *)
Hypothesis S3 :
  forall k, nrm (du k) <= (2*nu*Cb / hK k) * sqrt (aK k) * nrm (gu k).
\end{lstlisting}

\noindent
the theorem concludes
\begin{lstlisting}
Theorem abstract_stability :
  0 < C_stab c1 Cb xi C2  /\  BS >= C_stab c1 Cb xi C2 * NormSq.
\end{lstlisting}

\noindent
which is exactly $B_S(U_h,U_h)\ge \Cstab\,|||U_h|||^2$ with the explicit
$\Cstab=\min\{C_{\mathrm{visc}},C_u,1-C_2,1\}>0$. Everything between the
hypotheses and the conclusion --- the difference-of-squares cancellation, the
expansion of $\lVert G(u)\rVert^2$, the elemental Cauchy--Schwarz, the Young
step with parameter $\xi\bar C_{\mathrm{inv}}^2\nu\aK/h^2$, the exact
cancellation of the inverse-estimate constant, the coefficient collection, and
the summation over the mesh --- is machine-checked.

\subsection{What a green check does and does not assert}

An entry marked as covered means: \emph{the displayed relation is a theorem of
real arithmetic under the stated hypotheses}. For the algebraic files the
hypotheses are only positivity, so the theorem is unconditional in the
mathematics. For the two abstract lemmas the hypotheses additionally include
the named analytic base of Section~\ref{sec:left}; the green check asserts the
\emph{implication}, and the honesty of the result rests on that base being
exactly the standard, hand-audited facts it is labelled as.

% =====================================================================
\section{The crosswalk: theory \texorpdfstring{$\to$}{->} Coq}\label{sec:crosswalk}
% =====================================================================

Dependency and topic map of the twelve files. \file{Limits.v}/\file{DerivKit.v}
serve the concrete analytic files; \file{AbstractSums.v}/\file{InnerSpace.v}
serve the two abstract files, which also import the closed algebraic theorems
of \file{StabilityAlgebra.v}/\file{ContinuityAlgebra.v} and apply them per
element.

\begin{center}
\small
\begin{tabular}{@{}lll@{}}
\toprule
Coq file & Paper location & Role \\
\midrule
\file{Limits.v}      & --- & infrastructure: $\varepsilon$--$\delta$ limit theory \\
\file{DerivKit.v}    & --- & infrastructure: derivative toolkit \\
\file{AbstractSums.v}  & --- & infrastructure: finite mesh/face sums, discrete C--S \\
\file{InnerSpace.v}    & --- & infrastructure: abstract pre-Hilbert space \\
\file{StabilityAlgebra.v}  & \S\ref{sec:crossStab} (\S5)     & stability estimate, scalar algebra \\
\file{ContinuityAlgebra.v} & \S\ref{sec:crossCont} (App.\ B) & continuity proof, scalar algebra \\
\file{AbstractStability.v} & \S\ref{sec:crossAbsStab} (\S5) & \lbl{lemma:Stability}, complete theorem \\
\file{AbstractContinuity.v}& \S\ref{sec:crossAbsCont} (App.\ B) & \lbl{lemma:Continuity}, complete theorem \\
\file{TauDesign.v}   & \S\ref{sec:crossTau} (\S4 + App.\ C) & Fourier design of $\tau_1,\tau_2$ \\
\file{Asymptotics.v} & \S\ref{sec:crossAsym} (\S6)     & robustness in $(\Reh,\Dah)$ \\
\file{ManufacturedSolution.v} & \S\ref{sec:crossMMS} (\S7) & manufactured solution \\
\file{PlateauBump.v} & \S\ref{sec:crossBump} (\S7)     & plateau bump, through $C^1$ \\
\bottomrule
\end{tabular}
\end{center}

\subsection{Infrastructure: \texttt{InnerSpace.v} and \texttt{AbstractSums.v}}\label{sec:crossInfra}

\begin{longtable}{@{}L{0.42\textwidth}L{0.53\textwidth}@{}}
\toprule
\textbf{Coq statement} & \textbf{Content (all \emph{derived}, not assumed)} \\
\midrule \endhead
\multicolumn{2}{@{}l}{\emph{\file{InnerSpace.v} --- the elemental $L^2(K)$ pairing as a pre-Hilbert space}}\\
\coq{PreHilbert} (\coq{ip_sym}, \coq{ip_add_l}, \coq{ip_scal_l}, \coq{ip_pos}) & the four axioms modelling a symmetric positive-semidefinite $L^2(K)$ form. \\
\coq{ip_add_r}, \coq{ip_scal_r}, \coq{ip_sub_l/r}, \coq{ip_opp_l/r} & full bilinearity, from left-linearity $+$ symmetry. \\
\coq{CS_sq}, \coq{CS}, \coq{CS_le}, \coq{CS_ge} & Cauchy--Schwarz $|\langle x,y\rangle|\le\lVert x\rVert\lVert y\rVert$ (discriminant argument). \\
\coq{nrm_sq}, \coq{ip_expand_add}, \coq{ip_expand_sub} & Gram expansions $\lVert x\pm y\rVert^2$. \\
\coq{nrm_triangle} & the triangle inequality. \\
\coq{young_scalar}, \coq{young_vector} & the parametrized Young inequality at vector level. \\
\coq{diff_of_squares} & $\langle -a+b,\,a+b\rangle=\lVert b\rVert^2-\lVert a\rVert^2$: the ASGS cancellation. \\
\midrule
\multicolumn{2}{@{}l}{\emph{\file{AbstractSums.v} --- sums over the element / face list}}\\
\coq{Rsum}, \coq{Rsum_plus/minus/opp/scal}, \coq{Rsum_le}, \coq{Rsum_nonneg} & linearity and monotonicity of $\sum_{k\in\mathrm{Th}}$. \\
\coq{Rsum_abs_le} & the triangle inequality for sums. \\
\coq{Rsum_CS_sq}, \coq{Rsum_CS} & the discrete Cauchy--Schwarz over the mesh, by induction (the paper's ``C--S for the sums over elements''). \\
\coq{Rsum_sq_mono}, \coq{sqrt_plus_le}, \coq{sqrt_mono} & squared-sum monotonicity; subadditivity/monotonicity of $\sqrt{\cdot}$. \\
\bottomrule
\end{longtable}

\subsection{\texttt{StabilityAlgebra.v} --- \S5: the stability estimate (scalar algebra)}\label{sec:crossStab}

Section hypotheses: $\nu,h,\aK,\sigma>0$, $|\mathbf a|\ge 0$,
$c_1,c_2,\Cinv,\xi>0$. Mirrors \file{stability_estimate_verification.py}.

\begin{longtable}{@{}L{0.42\textwidth}L{0.53\textwidth}@{}}
\toprule
\textbf{Coq statement} & \textbf{Paper content} \\
\midrule \endhead
\coq{sigt_form_phi}, \coq{sigt_form_tau1}, \coq{sigt_form_key} & The three forms of $\sigt$ (\lbl{eq:SigmaAlpha}). \\
\coq{young_perfect_square}, \coq{young_inequality} & Young's step $-2xy\ge -x^2/\xi-\xi y^2$. \\
\coq{viscous_coefficient_expansion} & Viscous coefficient of \lbl{eq:StabilityEstimateFinal} $=$ the expanded \lbl{eq:ViscousCoefficientBound}. \\
\coq{velocity_coefficient_expansion}, \coq{velocity_slack_identity}, \coq{velocity_coefficient_lower_bound}, \coq{u_final_lower_bound} & \lbl{eq:VelocityCoefficientBound} and its \emph{exact} slack over $C_u\sigt$, $C_u=1-\xi\Cinv^2/c_1$. \\
\coq{eps_max_tau2_identity}, \coq{eps_tau2_le_C2}, \coq{pressure_term_coercive}, \coq{epsilon_smallness_chain} & The $\varepsilon$-smallness chain (\lbl{eq:UpperBoundOnEpsilon}): $\varepsilon\tau_2\le C_2$, pressure coeff $\ge(1-C_2)\varepsilon$. \\
\coq{viscous_coefficient_lower_bound}, \coq{visc_final_lower_bound} & Viscous bracket $\ge C_{\mathrm{visc}}\nu$, $C_{\mathrm{visc}}=\min\{2-4\Cinv^2/c_1,2(1-2/\xi)\}$. \\
\coq{stability_constants_positive} & $c_1>2\xi\Cinv^2$, $\xi>2$ $\Rightarrow$ coefficients positive (\lbl{eq:conditions_on_num_param}). \\
\coq{elemental_coercivity} & Elemental assembly $\ge\Cstab\times$ the triple-norm contributions, $\Cstab=\min\{C_{\mathrm{visc}},C_u,1-C_2,1\}$. \\
\bottomrule
\end{longtable}

\subsection{\texttt{ContinuityAlgebra.v} --- Appendix B: the continuity proof (scalar algebra)}\label{sec:crossCont}

Section hypotheses: $\nu,h,\aK,c_1,c_2>0$, $\aK\le 1$, $\sigma\ge 0$,
$|\mathbf a|\ge 0$. Explicit constants wherever the appendix writes a generic
$C$.

\begin{longtable}{@{}L{0.44\textwidth}L{0.51\textwidth}@{}}
\toprule
\textbf{Coq statement} & \textbf{Appendix content} \\
\midrule \endhead
\coq{P1_*}, \coq{P2_*}, \coq{P3_*}, \coq{P4_*}, \coq{P5_*} & \lbl{lem:parameters} (P1)--(P5) and the square-root forms. \\
\coq{sigt_id_tau}, \coq{sigt_id_sub}, \coq{one_minus_sigma_tau1} & \lbl{eq:sigmatilde} and the workhorse $1-\sigma\tau_1=\varphi_1\tau_1$. \\
\coq{keyvisc_sq}, \coq{keyvisc_sqrt}, \coq{T8_chain}, \coq{T13_chain} & \lbl{eq:keyvisc} (Step~4); the $\sigma\varphi_1\tau_1=\sigt$ absorptions (\lbl{eq:T8bound}, \lbl{eq:T13conv}). \\
\coq{volpart_scalar}, \coq{sqrtphi_tau_le}, \coq{phi_sqrttau_le}, \coq{sqrtphi_sqrttau_le_1} & \lbl{eq:volpart} (Step~6b) and Steps~6c--6d. \\
\coq{absorb1_*}, \coq{absorb2_*}, \coq{absorb4_*} & the \lbl{eq:absorb1}--\lbl{eq:absorb4} coefficient chains (Step~9). \\
\coq{jump_formula}, \coq{phi_diff_bound}, \coq{tau_comp_AB/_BA}, \coq{sg_comp_A_le_B}, \coq{jump_bound_A}, \coq{jumpsplit_AB}, \coq{III_absorb_AB/_BA} & \lbl{lem:jump}, \lbl{eq:jumpest}, \lbl{eq:jumpsplit}, Step~6d --- constants explicit ($C_J/c_J$). \\
\coq{delta_alpha_bound}, \coq{winv_ratio} & \lbl{H:porosity}'s derived bound and the \lbl{lem:winv} coefficient step $\alpha_\infty/\alpha_0^{1/2}\le\delta^{1/2}\alpha_\infty^{1/2}$. \\
\coq{lagrange2/3}, \coq{CS2}, \coq{CS3}, \coq{step1_bound}, \coq{norm5_absorption} & the finite Cauchy--Schwarz of Steps~1,~8 and the Step~9 squaring-and-summing. \\
\coq{tau2impl_le_tau2}, \coq{tau2_le_scaled_impl} & analysis-$\tau_2$ (\lbl{eq:Tau2Final}) vs implemented $\tau_2$ (\lbl{eq:Tau2}): equivalent up to $1+C_2$ (amendment~F6, \S\ref{sec:amend}). \\
\bottomrule
\end{longtable}

\subsection{\texttt{AbstractStability.v} --- \lbl{lemma:Stability} as a complete theorem}\label{sec:crossAbsStab}

Models each field's elemental restriction as a \coq{PreHilbert} vector and the
parameters via the closed \file{StabilityAlgebra.v} formulas. Conclusion:
\coq{abstract_stability} (see \S\ref{sec:absexample}). Trusted base:
\coq{HBS}, \coq{S3}, \coq{Heps} (the three of Section~\ref{sec:left}).

\begin{longtable}{@{}L{0.40\textwidth}L{0.55\textwidth}@{}}
\toprule
\textbf{Coq statement} & \textbf{Paper content (all machine-checked)} \\
\midrule \endhead
\coq{momentum_pairing}, \coq{mass_pairing} & the adjoint / difference-of-squares expansions of the residual pairings (\lbl{eq:StabilityEstimate}). \\
\coq{Bv_expand}, \coq{du_sq_bound} & expansion of $\lVert G(u)\rVert^2$; the Young step with the exact cancellation of the inverse-estimate constant $\bar C_{\mathrm{inv}}$ (\lbl{eq:BoundOfCrossedTerm}). \\
\coq{element_lower_bound}, \coq{element_coercive} & the per-element coercivity, importing \coq{elemental_coercivity}. \\
\coq{BS_as_sum} & re-summation of the tested form over the mesh. \\
\coq{abstract_stability} & \textbf{\lbl{lemma:Stability}}: $0<\Cstab$ and $B_S(U_h,U_h)\ge\Cstab\,|||U_h|||^2$. \\
\bottomrule
\end{longtable}

\subsection{\texttt{AbstractContinuity.v} --- \lbl{lemma:Continuity} as a complete theorem}\label{sec:crossAbsCont}

Adds a face list with the two neighbouring elements $e_1,e_2$ per interior
face. Conclusion: \coq{abstract_continuity}. Trusted base: the IBP / facewise
/ face-estimate / jump / mesh / inverse-estimate / $\varepsilon$ hypotheses of
Section~\ref{sec:left}.

\begin{longtable}{@{}L{0.40\textwidth}L{0.55\textwidth}@{}}
\toprule
\textbf{Coq statement} & \textbf{Appendix content (all machine-checked)} \\
\midrule \endhead
\coq{BS} $= $ \coq{T1} $+\cdots+$ \coq{T18} & the eighteen-term decomposition of $B_S$ (\lbl{eq:Bstab}, Step~0). \\
\coq{bound_T1}, \coq{bound_T3_T14}, \coq{bound_T5}, \coq{bound_T10}, \coq{bound_T15}, \coq{bound_T16}, \coq{bound_T17}, \coq{bound_T18}, \coq{keyvisc_coef} & the per-term bounds of Steps~1--5 and~7. \\
\coq{step6a} & the exact five-term rewriting of Step~6a (\lbl{eq:fiveterms}), consuming the IBP identities \lbl{eq:elemibp} etc. \\
\coq{bound_VolP}, \coq{bound_JmpP} (\coq{Jb11}--\coq{Jb22}) & Step~6b: volumetric part via $\varphi_1\tau_1\le\sqrt{\varphi_1}\sqrt{\tau_1}$ and (P3); jump part via \lbl{lem:jump} through four face-resolved bounds, facewise C--S, bounded multiplicity. \\
\coq{bound_II_V} & Step~6c. \\
\coq{bound_III} (\coq{JD11}--\coq{JD22}) & Step~6d, with the four face-difference bounds. \\
\coq{face_CS_pair} & facewise Cauchy--Schwarz with the multiplicity hypotheses. \\
\coq{abstract_continuity_sharp} & Step~8: the sharp porosity-weighted assembly (\lbl{eq:assembly}). \\
\coq{absorb_elem}, \coq{step9} & Step~9: the \lbl{eq:absorb1}--normconv absorption chain. \\
\coq{abstract_continuity} & \textbf{\lbl{lemma:Continuity}}: $B_S(U,V)\le C_{\mathrm{tot}}(\mathrm{Br}_U+\mathrm{Br}_P)\,|||V|||$, $C_{\mathrm{tot}}$ fully explicit. \\
\bottomrule
\end{longtable}

\subsection{\texttt{TauDesign.v} --- \S4 / Appendix C: Fourier design of \texorpdfstring{$\tau$}{tau}}\label{sec:crossTau}

Mirrors \file{fourier_tau_verification.py}. Section hypotheses: $\nu,\alpha,h>0$.

\begin{longtable}{@{}L{0.40\textwidth}L{0.55\textwidth}@{}}
\toprule
\textbf{Coq statement} & \textbf{Paper content} \\
\midrule \endhead
\coq{viscous_symbol_closed_form} & $\sum_{ij}k_ik_j\mathbf K_{ij}=\alpha\nu(|k|^2 I+\tfrac13 kk^\top)$, $d{=}3$ (\lbl{eq:matrices_stationary_strong_problem}). \\
\coq{viscous_parallel_eigenpair}, \coq{viscous_transverse_eigenpair}, \coq{parallel_factor_d3/d2} & viscous spectrum $\tfrac43\alpha\nu|k|^2$, $\alpha\nu|k|^2$; the $d{=}3,2$ prefactors. \\
\coq{convective_symbol_scalar}, \coq{convective_spectral_radius} & convective symbol and radius $(\alpha/h)|w\cdot k|$. \\
\coq{coupling_*}, \coq{coupling_spectrum}, \coq{coupling_spectral_radius} & pressure--velocity coupling $\mathbf C$: spectrum $\{0,\pm|k|\}$. \\
\coq{taub_pair_solves}, \coq{taub_pair_unique} & $XY=c^2,\;X/Y=\lambda$ $\Rightarrow$ unique positive $(c\sqrt\lambda,c/\sqrt\lambda)$. \\
\coq{sqrt_lam_choice}, \coq{tau1_inv_assembled}, \coq{tau2_assembled}, \coq{visc_plus_conv_momentum} & assembly of \lbl{eq:Tau1}, \lbl{eq:Tau2}, \lbl{eq:CAlpha}, \lbl{eq:TauNavierStokes}. \\
\bottomrule
\end{longtable}

\subsection{\texttt{Asymptotics.v} --- \S6: robustness in \texorpdfstring{$(\Reh,\Dah)$}{(Re,Da)}}\label{sec:crossAsym}

Mirrors \file{robustness_asymptotics_verification.py}. Section hypotheses:
$\nu,h,\aK>0$; $\Reh,\Dah\ge0$.

\begin{longtable}{@{}L{0.42\textwidth}L{0.53\textwidth}@{}}
\toprule
\textbf{Coq statement} & \textbf{Paper content} \\
\midrule \endhead
\coq{tau1_general_form}, \coq{tau2_general_form}, \coq{sigt_general_form} & closed forms in $(\Reh,\Dah)$ (\lbl{eq:GeneralAsymptoticBehaviourOfParameters}). \\
\coq{dom_viscosity_*} & dominant viscous diffusion ($\Reh,\Dah\to0$) (\S6.1). \\
\coq{dom_convection_*} & dominant convection ($\Reh\to\infty$): genuine \coq{tendsto_at_top} limits (\S6.2). \\
\coq{dom_reaction_*} & dominant reaction ($\Dah\to\infty$): $\tau_1\sigma\to1$, $\Dah$-independence of $\tau_2$ (paper \S6.3). \\
\coq{convection_coefficient}, \coq{viscous_coefficient}, \coq{pressure_coefficient}, \coq{reaction_coefficient}, \coq{forcing_coefficient} & nondimensionalisation coefficients $\Reh,1,(1{+}\Reh{+}\Dah),\Dah$ (\lbl{eq:DimensionlessMomentumEquation}, \lbl{eq:PressureScale}). \\
\bottomrule
\end{longtable}

\subsection{\texttt{ManufacturedSolution.v} --- \S7: the manufactured solution}\label{sec:crossMMS}

Mirrors \file{manufactured_solution_verification.py}; genuine
\coq{derivable_pt_lim} derivatives.

\begin{longtable}{@{}L{0.42\textwidth}L{0.53\textwidth}@{}}
\toprule
\textbf{Coq statement} & \textbf{Paper content} \\
\midrule \endhead
\coq{manufactured_mass_conservation_2d/3d} & $\nabla\!\cdot(\alpha\mathbf u)=0$ for arbitrary differentiable positive $\alpha$ (\lbl{eq:ManufacturedProblem}). \\
\coq{boundary_trace_nonzero} & $u_2(0,0)>0$: the trace does not vanish (amendment A11). \\
\coq{rhs_closed_form}, \coq{denominator_gap}, \coq{mid_le_rhs}, \coq{eps_choice_margin} & the \lbl{eq:EpsilonRef} chain at $C_2=10^{-2}$, and $\varepsilon:=10^{-4}\varepsilon_{\mathrm{ref}}$. \\
\coq{dbf_a_nonneg}, \coq{dbf_b_nonneg} & DBF coefficients admissible on $0<\alpha\le1$ (\lbl{eq:DBFResistanceTerm}). \\
\coq{alpha_lit_at_1}, \coq{alpha_lit_deriv(_pos)}, \coq{alpha_lit_at_0_range} & literature profile $\alpha(y)=0.45+0.55e^{y-1}$ (paper \S7.3). \\
\bottomrule
\end{longtable}

\subsection{\texttt{PlateauBump.v} --- \S7: the plateau bump, through \texorpdfstring{$C^1$}{C1}}\label{sec:crossBump}

$\gamma(\eta)=(2\eta-1)/(\eta(1-\eta))$,
$\alpha_{\mathrm{mid}}=1-(1-\alpha_0)/(1+e^\gamma)$
(\lbl{eq:PlateauBumpFunction}, \lbl{eq:Gamma}). Explicit $\varepsilon$--$\delta$
on $t<e^t$ and $t^3/27\le e^t$.

\begin{longtable}{@{}L{0.40\textwidth}L{0.55\textwidth}@{}}
\toprule
\textbf{Coq statement} & \textbf{Paper content} \\
\midrule \endhead
\coq{gam_deriv}, \coq{gam'_pos}, \coq{q_deriv}, \coq{am_deriv}, \coq{am'_pos} & genuine derivatives of $\gamma$ and $\alpha_{\mathrm{mid}}$; strict monotonicity. \\
\coq{am_range} & $\alpha_0<\alpha_{\mathrm{mid}}<1$ unconditionally. \\
\coq{exp_gam_small}, \coq{exp_gam_cubic_small}, \coq{gam'_upper_small/_big} (and mirrors) & quantitative decay bounds near $0$ and $1$. \\
\coq{am_limit_0}, \coq{am_limit_1} & continuity of the joins ($C^0$). \\
\coq{am'_limit_0}, \coq{am'_limit_1} & vanishing derivative at the joins ($C^1$). \\
\bottomrule
\end{longtable}

The paper's $C^\infty$ claim rests on every derivative vanishing at the joins;
the \emph{first-order} case (which the SymPy script also checks) is what is
formalised.

\subsection{Infrastructure: \texttt{Limits.v} and \texttt{DerivKit.v}}

\file{Limits.v} defines \coq{tendsto_at_top}, \coq{tendsto_at_0plus},
\coq{tendsto_at_1minus} (explicit $\varepsilon$--$\delta$, no filters), transfer
lemmas, and the generic rational limits consumed by \file{Asymptotics.v} and
\file{PlateauBump.v}. \file{DerivKit.v} provides \coq{derivable_pt_lim}
helpers and derivatives of $a\cdot t$, $\sin(at)$, $\cos(at)$, $c\cdot f(t)$,
$e^{t-1}$. Neither is paper-specific.

% =====================================================================
\section{The residual trusted base and what is left}\label{sec:left}
% =====================================================================

With the abstract layer in place, the ``what is left'' story is now a short,
explicit list. Every item below is either (i) a named hypothesis of the two
abstract theorems --- standard analysis, hand-audited and slated for Lean --- or
(ii) genuinely outside the current scope, and covered by the SymPy suite.

\subsection{The named trusted base of the two abstract theorems}

The companion \file{LEAN_ROADMAP.md} classifies each hypothesis into three
classes. Class~D items are assumptions on data/mesh --- hypotheses of the
paper's theorems in any formal system, with nothing to prove. Classes~G and~I
are the textbook analytic facts (Brenner--Scott, Ern--Guermond), hand-audited
against the manuscript in \file{AUDIT.md}~\S4.

\begin{longtable}{@{}L{0.30\textwidth}L{0.34\textwidth}cL{0.16\textwidth}@{}}
\toprule
\textbf{Coq hypothesis} & \textbf{Paper} & \textbf{Class} & \textbf{Status} \\
\midrule \endhead
\coq{HBS} & Galerkin/adjoint identity (\lbl{eq:StabilityEstimate} testing) & G & hand-audited \\
\coq{S3} & \lbl{eq:winv-divvisc} (\lbl{lem:winv}) & I & hand-audited \\
\coq{Heps}, \coq{H_eps} & \lbl{eq:UpperBoundOnEpsilon} / \lbl{eq:epscond} & D & assumption \\
\coq{H_skew} & \lbl{eq:skew}, identity (i) & G & hand-audited \\
\coq{H_ibp_vp}, \coq{H_ibp_qu} & \lbl{eq:globalibp}, identity (ii) & G & hand-audited \\
\coq{H_elem_conv_ibp} & \lbl{eq:elemibp}, identity (iii) & G & hand-audited \\
\coq{H_elem_p_ibp} & elemental pressure IBP (Step~6b) & G & hand-audited \\
\coq{H_assemble_p}, \coq{H_assemble_c} & facewise assembly (Steps~6b,~6d) & G & hand-audited \\
\coq{H_face_p}, \coq{H_face_c} & \lbl{eq:jumpest}/face estimates: H\"older $+$ $\mathrm{meas}\le Ch^{d-1}$ $+$ $L^\infty$ inverse $+$ $\alpha\le\aK$ & I,G & hand-audited \\
\coq{H_jump} & \lbl{H:jump} & D & assumption \\
\coq{H_mult1}, \coq{H_mult2} & bounded face multiplicity (\lbl{H:mesh}) & D & assumption \\
\coq{Hw_gu} \dots \coq{Hw_divu} (8) & weighted inverse estimates of \lbl{lem:winv} (\lbl{eq:inverse}) & I & hand-audited \\
\bottomrule
\end{longtable}

\noindent
The roadmap observes that Classes~G and~I collapse onto two primitives:
\begin{itemize}[nosep]
  \item \textbf{Class~I} reduces to the \emph{unweighted} inverse estimate on a
    shape-regular simplex (an $L^2$ and an $L^\infty$ version) plus a
    weighted-to-unweighted reduction --- whose \emph{scalar} chain
    ($\alpha_\infty/\alpha_0$ comparabilities under \lbl{H:porosity}) is
    \emph{already machine-checked} in \file{ContinuityAlgebra.v}
    (\coq{delta_alpha_bound}, \coq{winv_ratio}); only the norm-level gluing
    remains.
  \item \textbf{Class~G} reduces to the divergence theorem on a single simplex
    for (polynomial)$\times$($W^{1,\infty}$ weight) data, plus finite
    bookkeeping over elements and faces using continuity of the integrands.
\end{itemize}

\subsection{Hand-audited without a Coq statement}

The interpolation-theoretic replacements of \lbl{lem:continterp} (and the
triangle-inequality assembly of \lbl{thm:convergence} / \lbl{th:Convergence}
built on them) are re-derived line by line against the manuscript in
\file{AUDIT.md}~\S4 but do not yet have a Coq statement. This is the one part
of the stability/continuity development still resting on the hand audit alone.

\subsection{The Lean~4 roadmap (a plan, not a verification)}

\file{LEAN_ROADMAP.md} lays out an assistant-driven plan to discharge
Classes~G and~I in Lean~4 + Mathlib, with a companion statement skeleton
\file{PorousNSToolbox.lean}. \textbf{That skeleton is explicitly a statement
roadmap, not a proof artifact}: every target ends in \coq{sorry} and it is not
expected to elaborate against current Mathlib without API repair. Nothing in
the Lean file is machine-checked; it pins the intended semantics of each
residual target (unweighted inverse estimate on a simplex, its $L^\infty$
sibling and the face-measure bound, the weighted reductions, and the
simplex divergence theorem) for a future formalisation loop. The dependency-minimal
path noted there is the inverse-estimate technology, which alone would
discharge \coq{S3} and all eight \coq{Hw_*}, shrinking
\lbl{lemma:Stability}'s trusted base to the single Green identity \coq{HBS}.

\subsection{Deliberately delegated to the SymPy suite}

The Appendix~A elemental matrices
(\file{elemental_matrices_verification.py},
\file{elemental_bilinear_form_verification.py}) are transcription checks by
symbolic differentiation of shape functions --- a computer-algebra task where a
proof assistant adds cost but no assurance. They are left to SymPy.

\subsection{Out of scope for standard-library Coq}

Two items need machinery absent from stdlib Coq (natural in Lean~4 + Mathlib),
covered only by the SymPy scripts: the subscale-norm equivalence of
\file{subscale_norm_verification.py} (Appendix~B9, Rayleigh-quotient spectral
theory), and the CDR adjoint / Green's-identity computation of
\file{cdr_operator_verification.py} (the divergence theorem in general form).

\subsection{The two standing caveats}

\begin{itemize}[nosep]
  \item \textbf{Statements, not models.} Coq certifies that each identity,
    inequality and limit follows from the real-number axioms (and, for the two
    abstract lemmas, the named base). It cannot certify that these are the
    \emph{right} statements about the physics --- that requires reading them,
    which the crosswalk is meant to make easy.
  \item \textbf{Theory, not code.} Coq says nothing about whether the Julia or
    Kratos solvers implement the verified formulas. That correspondence is
    established empirically by the method-of-manufactured-solutions convergence
    tests in the codebase (the $O(h^{k+1})$ plateaus).
\end{itemize}

% =====================================================================
\section{Theory-side amendments applied for compatibility}\label{sec:amend}
% =====================================================================

Formalising the proofs surfaced a few points where the manuscript's hypotheses
or notation were slightly looser than what the proofs consume. The following
amendments (catalogued in \file{proof_verification/AUDIT.md}, findings F1--F7)
have been \textbf{applied} to \file{theory/paper/} so the manuscript is
$100\%$ compatible with the verified statements. None changes a mathematical
result; each closes a precision or hygiene gap. The added text is marked with
the manuscript's \verb|\amend{}| review macro.

\begin{longtable}{@{}cL{0.72\textwidth}c@{}}
\toprule
\textbf{ID} & \textbf{Amendment} & \textbf{Status} \\
\midrule \endhead
F1 & \lbl{H:advection} now grants $\mathbf a|_K\in W^{1,\infty}(K)^d$ elementwise (with the localisation of $\nabla\!\cdot(\alpha\mathbf a)=0$ to each $K$), which identity (iii)/\lbl{eq:elemibp} --- the Coq hypothesis \coq{H_elem_conv_ibp} --- actually uses. & applied \\
F2 & \lbl{H:porosity}'s derived bound $\alpha_{\infty,K}\le\delta_\alpha\alpha_{0,K}$ now states the element-convexity it relies on (matching \coq{delta_alpha_bound}/\coq{winv_ratio}). & applied \\
F3 & the two previously-unlabelled coefficient displays after \lbl{eq:UpperBoundOnEpsilon} now carry stable labels \lbl{eq:ViscousCoefficientBound} and \lbl{eq:VelocityCoefficientBound} (the anchors of \coq{visc_final}/\coq{u_855}); the SymPy headers were updated to match. & applied \\
F4 & the main-text triple norm now writes the elementwise-weighted term $\lVert\sigt_\alpha^{1/2}\mathbf u_h\rVert_h^2$ with the broken-norm subscript, matching \lbl{eq:triplenorm}. & applied \\
F6 & a remark after \lbl{eq:Tau2Final} records that the analysis $\tau_2$ and the implemented \lbl{eq:Tau2} are equivalent up to $1+C_2<2$ (machine-checked: \coq{tau2impl_le_tau2}, \coq{tau2_le_scaled_impl}). & applied \\
F7 & a clause notes that the coefficient bounds are written for $\aK=\alpha_{\infty,K}$; other representatives cost factors of $\delta_\alpha$. & applied \\
\midrule
F5 & \file{supplement.tex} is untouched SIAM template boilerplate --- a submission-set decision for the authors, left unchanged. & deferred \\
\bottomrule
\end{longtable}

\noindent
After these edits the paper builds clean (\coq{latexmk} exit~0, $0$ undefined
references/citations); the two new equation labels resolve.

% =====================================================================
\section{Reproducing the check}\label{sec:repro}
% =====================================================================

From \file{proof_verification/coq-formal/}:
\begin{lstlisting}[language=bash,keywordstyle=\color{black}]
# one-time: install Coq / Rocq (any of)
brew install coq          # macOS (installs Rocq 9.1.x, provides coqc/coqchk)
apt install coq           # Debian/Ubuntu (Coq 8.x)
opam install coq

# verify
./run_all.sh              # per-file PASS + coqchk kernel re-verification
\end{lstlisting}

\noindent
All twelve files compile and kernel-check unchanged on Rocq~9.1.1 (the only
console output is a deprecation notice that \coq{From Coq} is now spelled
\coq{From Stdlib}; the $8.16$--$8.20$-compatible spelling is kept on purpose
for portability). A passing run prints one \coq{[PASS]} line per file, then
\emph{``coqchk: all modules kernel-verified''}.

\bigskip
\noindent\emph{Cross-references.} File-level detail and the SymPy$\leftrightarrow$Coq
mapping: \file{README.md} in the \file{coq-formal/} directory. The step-by-step
audit of the analytic steps and the F1--F7 findings: \file{AUDIT.md}. The Lean~4
plan for the residual base: \file{LEAN_ROADMAP.md} and \file{PorousNSToolbox.lean}.
The authoritative theory: \file{theory/paper/article.tex} and its appendices
(\file{continuity_appendix.tex}, \file{fourier_appendix.tex},
\file{elemental_matrices_appendix.tex}).

\end{document}
